% GENERATED FILE. Edit canonical lessons, then run npm run generate:books.
% canonical-source-hash: fnv1a64:ba6cad37a6debdfb
% canonical-lessons: FR-C199-entree, FR-C199-sortie, FR-C199-bagages, FR-C199-siege, FR-C199-volant

\chapter{Entrance, Exit, Luggage, Seat, Steering wheel}
\label{ch:fr-entrance-exit-luggage-seat-steering-wheel}
\hlchaptermodality{199}

\begin{chapteropening}
\textbf{By the end of this chapter:} \emph{I can say an entrance, an exit, luggage, a seat and a steering wheel.}

Complete the last lesson of chapter 199: I can say an entrance, an exit, luggage, a seat and a steering wheel.
\end{chapteropening}

\section[l'entrée]{l'entrée — an entrance}

\label{lesson:FR-C199-entree}



\begin{quote}
Before the new one: say the French for a coin, a room, then the French for a banknote.
\end{quote}

\subsection*{You'll want to know: l'entrée}
\textbf{l'entrée} — \textquotedblleft{}an entrance\textquotedblright{}.

\begin{cousinweb}
\textbf{entrée} is the past participle of the verb for \textquotedblleft{}to enter\textquotedblright{}, used as a noun, $\leftarrow$ Latin \emph{intrare} (Lewis \& Short): the way in.
\end{cousinweb}

\begin{grammarlens}[title={What you've built}]
One more word for this chapter.
\end{grammarlens}

\subsection*{Your turn}
\begin{itemize}
\raggedright
  \item \emph{Say it:} \emph{l'entrée}
  \item \emph{Say it:} \emph{l'entrée}, once more
  \item \emph{Say it:} say \emph{l'entrée}
  \item \emph{Recall:} say the French for a coin, a room, then the French for a banknote, then say \emph{l'entrée} again
\end{itemize}

\begin{tcolorbox}[breakable,colback=teal!4,colframe=teal!35!black,title={Before you move on}]
What does l'entrée mean? (\textquotedblleft{}An entrance\textquotedblright{}.) Say it once more.
\end{tcolorbox}
\section[la sortie]{la sortie — an exit}

\label{lesson:FR-C199-sortie}



\begin{quote}
Before the new one: say the French for a banknote, then the French for an entrance.
\end{quote}

\subsection*{You'll want to know: la sortie}
\textbf{la sortie} — \textquotedblleft{}an exit\textquotedblright{}.

\begin{cousinweb}
\textbf{sortie} is the past participle of the verb for \textquotedblleft{}to go out\textquotedblright{}, used as a noun: the way out. The verb's origin is uncertain; the usual guess is Latin \emph{sortiri}, \textquotedblleft{}to draw lots\textquotedblright{}.
\end{cousinweb}

\begin{grammarlens}[title={What you've built}]
One more word for this chapter.
\end{grammarlens}

\subsection*{Your turn}
\begin{itemize}
\raggedright
  \item \emph{Say it:} \emph{la sortie}
  \item \emph{Say it:} \emph{la sortie}, once more
  \item \emph{Say it:} say \emph{la sortie}
  \item \emph{Recall:} say the French for a banknote, then the French for an entrance, then say \emph{la sortie} again
\end{itemize}

\begin{tcolorbox}[breakable,colback=teal!4,colframe=teal!35!black,title={Before you move on}]
What does la sortie mean? (\textquotedblleft{}An exit\textquotedblright{}.) Say it once more.
\end{tcolorbox}
\section[les bagages]{les bagages — luggage}

\label{lesson:FR-C199-bagages}



\begin{quote}
Before the new one: say the French for an entrance, then the French for an exit.
\end{quote}

\subsection*{You'll want to know: les bagages}
\textbf{les bagages} — \textquotedblleft{}luggage\textquotedblright{}.

\begin{cousinweb}
\textbf{bagages} — built from Old French \emph{bague}, \textquotedblleft{}a bundle\textquotedblright{} (EtymDB 2.1), with \emph{-age}; \emph{bague}'s own origin is uncertain.
\end{cousinweb}

\begin{grammarlens}[title={What you've built}]
One more word for this chapter.
\end{grammarlens}

\subsection*{Your turn}
\begin{itemize}
\raggedright
  \item \emph{Say it:} \emph{les bagages}
  \item \emph{Say it:} \emph{les bagages}, once more
  \item \emph{Say it:} say \emph{les bagages}
  \item \emph{Recall:} say the French for an entrance, then the French for an exit, then say \emph{les bagages} again
\end{itemize}

\begin{tcolorbox}[breakable,colback=teal!4,colframe=teal!35!black,title={Before you move on}]
What does les bagages mean? (\textquotedblleft{}Luggage\textquotedblright{}.) Say it once more.
\end{tcolorbox}
\section[le siège]{le siège — a seat}

\label{lesson:FR-C199-siege}



\begin{quote}
Before the new one: say the French for an exit, then the French for luggage.
\end{quote}

\subsection*{You'll want to know: le siège}
\textbf{le siège} — \textquotedblleft{}a seat\textquotedblright{}.

\begin{cousinweb}
\textbf{siège} $\leftarrow$ Old French \emph{siege} $\leftarrow$ Vulgar Latin \emph{*sedicum}, from Latin \emph{sedes}, \textquotedblleft{}a seat\textquotedblright{} — inherited. (EtymDB 2.1)
\end{cousinweb}

\begin{grammarlens}[title={What you've built}]
One more word for this chapter.
\end{grammarlens}

\subsection*{Your turn}
\begin{itemize}
\raggedright
  \item \emph{Say it:} \emph{le siège}
  \item \emph{Say it:} \emph{le siège}, once more
  \item \emph{Say it:} say \emph{le siège}
  \item \emph{Recall:} say the French for an exit, then the French for luggage, then say \emph{le siège} again
\end{itemize}

\begin{tcolorbox}[breakable,colback=teal!4,colframe=teal!35!black,title={Before you move on}]
What does le siège mean? (\textquotedblleft{}A seat\textquotedblright{}.) Say it once more.
\end{tcolorbox}
\section[le volant]{le volant — a steering wheel}

\label{lesson:FR-C199-volant}



\begin{quote}
Before the new one: say the French for luggage, then the French for a seat.
\end{quote}

\subsection*{You'll want to know: le volant}
\textbf{le volant} — \textquotedblleft{}a steering wheel\textquotedblright{}.

\begin{cousinweb}
\textbf{volant} is the present participle of the verb for \textquotedblleft{}to fly\textquotedblright{}, $\leftarrow$ Latin \emph{volare} (Lewis \& Short), used as a noun: a thing that spins freely.
\end{cousinweb}

\begin{grammarlens}[title={What you've built}]
One more word for this chapter.
\end{grammarlens}

\subsection*{Your turn}
\begin{itemize}
\raggedright
  \item \emph{Say it:} \emph{le volant}
  \item \emph{Say it:} \emph{le volant}, once more
  \item \emph{Say it:} say \emph{le volant}
  \item \emph{Recall:} say the French for luggage, then the French for a seat, then say \emph{le volant} again
\end{itemize}

\begin{tcolorbox}[breakable,colback=teal!4,colframe=teal!35!black,title={Before you move on}]
What does le volant mean? (\textquotedblleft{}A steering wheel\textquotedblright{}.) Say it once more.
\end{tcolorbox}
